\songtitle{Dapich Sápica}

\tag*{1997}\tag{tango}\tag{spanish-lyrics}

\begin{notes}
A moody tango-infused piece that blends dramatic lyricism with tense rhythmic motion. %
The lyrics begin with “Hace dos años que partió el grande,”, giving a quick sense of the song’s tone.
\end{notes}

\begin{style}
tango, milonga, dramatic storytelling, 132 BPM, bandoneon stabs, nylon-string guitar rasgueo, upright bass pizzicato, piano tresillo figures, violin counterline, male tenor lead, spoken phrasing, theatre choir, tape saturation, vintage mono horn room, syncopated handclaps, minor key march, cinematic intro, bittersweet reveal
\end{style}

\begin{lyrics}
Hace dos años que partió el grande,\\
el sabio Dapich, orgullo de Sápica.\\
Los mismos necios que gritaron “¡que se vaya!”\\
hoy lloran su ausencia con el alma marchita.

Es invierno feroz, la nieve castiga,\\
los pinos se doblan bajo el peso del hielo.\\
El viento aúlla como un lobo en la cima\\
y el frío muerde más hondo que un acero.

Los hombres, con manos agrietadas y duras,\\
subidos a techos que crujen de pena,\\
clavan y remiendan con furia y premura\\
mientras el cielo amenaza tormenta.

Y en medio del blanco, un niño inocente,\\
jugando con nieve que tiñe de frío,\\
ve una figura que avanza lentamente—\\
negra contra el blanco, como un mal augurio.

Los viejos de Sápica ya saben quién viene,\\
las mujeres se asoman con ojos cansados.\\
Todos lo saben… pero nadie lo dice,\\
porque duele el regreso del que nada ha hallado.

Camina vencido, el poncho hecho harapos,\\
espaldas cargadas con toda la ausencia.\\
El niño lo mira, le tiembla la mano…\\
y aún así le sonríe con ternura inmensa.

Pasa a su lado, lento, sin palabras.\\
El viento se calla, parece que espera.\\
Y ahí, frente a toda la aldea callada,\\
cae la verdad como un puñal en la niebla—

No era el sabio.

Era su hijo.

\end{lyrics}

